% !TEX root = scientific-computing.tex
% !TEX encoding = UTF-8 Unicode

\chapter{Probability of Poker Hands} \label{ch:poker}

This chapter will examine coming up with poker hands. If you are unfamiliar with poker or the hands of poker, here's a quick synopsis and the \href{https://en.wikipedia.org/wiki/List_of_poker_hands}{Wikipedia page on Poker Hands}.  And recall from Chapter \ref{ch:comp-types} that a card has both a rank (Ace, 2--10, Jack, Queen, King) and a suit (hearts, diamonds, clubs and spades). 

In this chapter, we are only concerned with 5 cards with no jokers and just determining if a hand satisfies one of the following:

\begin{itemize}
\item \emph{Royal Flush:} the ranks of the cards are 10, J, Q, K, A and all cards have the same suit. 
\item \emph{Straight Flush:} the ranks are sequential and all cards have the same suit. 
\item \emph{Flush:} All cards have the same suit.  Generally excludes straight flushes. 
\item \emph{Straigh:} The ranks of the 5 cards are sequential.  Generally exclude straight flushes. 
\item \emph{Four of a kind:} four of the cards have the same rank
\item \emph{Full House:} two cards have the same rank, the other three cards have the same rank.  The suit doesn't matter. 
\item \emph{Three of a kind:} three of the cards have the same rank.  The other two cards do not have the same rank.  That is, don't count full houses. 
\item \emph{Two pairs:} Two cards have the same rank.  Two of the remaining cards have the same rank, but different than the first two pair.  The 5th card does not make it a full house. 
\item \emph{One pair:} two cards have the same rank.  The remaining cards do not make it a different type of hand (full house, three of a kind, etc.) 
\item \emph{No pair or nothing:} the cards don't form any other hand. 
\end{itemize}

\section{A user-defined package}

In Chapter \ref{ch:modules}, we will create a module, but it is helpful for the topics in this chapter to use that module.  Download \texttt{PlayingCards.jl} from ???.  This module contains the \verb!Card! and \verb!Hand! types we developed in Chapter \ref{ch:comp-types}.  

This is a module/package, which like other packages, need to be loaded	
with either the \texttt{using} or \texttt{import} keyword. Since this is
just a file in the current directory, we first run the file and then
load it

\begin{jlverbatim}
include("PlayingCards.jl")
using .PlayingCards
\end{jlverbatim}
\begin{jlcode}[poker]
include("../julia-files/PlayingCards.jl")
using .PlayingCards
\end{jlcode}

%
where the \texttt{.} represents a local (current directory) module.  Note that when running a module that has been added (downloaded), no . is needed.  

% note: we have loading in this module in the preamble of the text: 


Once we have the function written, we should test it on a few known and
unknown full house hands. Try testing:
%
\begin{jlblock}[poker]
fh1 = Hand([Card(4,1),Card(4,3),Card(4,4),Card(7,1),Card(7,2)])
fh2 = Hand([Card(4,1),Card(4,3),Card(7,4),Card(7,1),Card(7,2)])
fh3 = Hand([Card(2,1),Card(4,3),Card(4,4),Card(7,1),Card(7,2)])
\end{jlblock}
%
The first 2 are full house hands and the last is not.


Now, let's perform a Monte Carlo simulation on a large number of poker
hands and test if this gives the result we want.  We'd like to have a function that tests many different hands.  In general, this is called a trial, so we'll have following which passes in a function that takes a hand and a number of trials and returns the fraction of times that hand satisfies that function.  First, here's the \verb!runTrial! function. 

\begin{jlblock}[poker]
using Random
function runTrial(f::Function, trials::Integer)
  deck=collect(1:52) # creates the array [1,2,3,...,52]
  numhands=0
  for i=1:trials
    shuffle!(deck)
    h = Hand(map(Card,deck[1:5]))  # creates a hand of the first five cards of the shuffled deck
    if(f(h))
        numhands+=1
    end
  end
  numhands/trials
end
\end{jlblock}


%
This returns \jlc[poker]{print(runTrial(isFullHouse,10_000))} \printpythontex[verb]. 
It would be helpful to put this in a function. We will do this and put all
of these in a package together in the chapter \ref{ch:modules}. 
